\documentclass[11pt,a4paper]{article}

\usepackage[utf8]{inputenc}
\usepackage[T1]{fontenc}
\usepackage{amsmath,amssymb}
\usepackage{booktabs}
\usepackage{hyperref}
\usepackage{xcolor}
\usepackage{colortbl}
\usepackage[margin=2.5cm]{geometry}
\usepackage{caption}
\usepackage{subcaption}
\usepackage{array}
\usepackage{listings}

\definecolor{bestcell}{rgb}{0.85,1.0,0.85}
\definecolor{badcell}{rgb}{1.0,0.88,0.85}
\definecolor{tlcell}{rgb}{0.95,0.95,0.95}
\definecolor{headercolor}{rgb}{0.22,0.39,0.60}

\lstset{basicstyle=\ttfamily\small, breaklines=true, frame=single}

\hypersetup{
  colorlinks=true,
  linkcolor=headercolor,
  urlcolor=headercolor,
  citecolor=headercolor
}

\title{\textbf{Barrier vs.\ Simplex for LP Relaxations in Clp}\\[0.4em]
       \large Comparing Native Cholesky, AMD Cholesky, Primal Simplex,\\
              and Dual Simplex on MIPLIB~2017 Instances}
\author{Experimental study --- CBC/CLP development workspace}
\date{March 2026}

% ─────────────────────────────────────────────────────────────────────────────
\begin{document}
\maketitle
\tableofcontents
\bigskip

% ─────────────────────────────────────────────────────────────────────────────
\section{Motivation}

The \texttt{-guess} option in Clp (documented separately in
\texttt{doc/guess\_option.pdf}) automatically selects an LP algorithm based on
instance features.  In its current form it always chooses either primal or dual
simplex; the barrier (interior-point) branch is present in the code but is dead.
A natural question is: \emph{is that correct?}  When, if ever, does barrier
outperform simplex on the LP relaxations that arise inside a MIP solve?

A secondary question is whether the SuiteSparse / CHOLMOD AMD reordering
(available in this build via \texttt{-chol~uni}) improves the barrier solver
relative to Clp's native Cholesky, which is relevant both for speed and for
peak memory consumption.

This report documents a systematic experiment that benchmarks four solver
configurations on eleven MIPLIB~2017 instances that are cited in the
\texttt{guess\_option} document as structural candidates for the barrier
method.

% ─────────────────────────────────────────────────────────────────────────────
\section{Experimental Setup}

\subsection{Solver}

All runs use the \texttt{clp} binary from the local development build
(\texttt{\$HOME/prog/cbc/bin/clp}, version \emph{devel}, build date
2026-03-24), compiled with:
\begin{itemize}
  \item GCC~15 (\texttt{-O3 -march=native});
  \item SuiteSparse / AMD (\texttt{libamd}, \texttt{libcholmod}) from the
        system package;
  \item OpenBLAS for dense linear algebra.
\end{itemize}

\subsection{Solver Configurations}

Four Clp configurations are tested, differing only in LP algorithm and
Cholesky ordering:

\begin{table}[h]
\centering
\caption{Solver configurations.}
\label{tab:configs}
\begin{tabular}{@{}lll@{}}
\toprule
\textbf{Label} & \textbf{Clp command-line flags} & \textbf{Description} \\
\midrule
\texttt{barrier\_default}  & \texttt{-seconds 28800 -barrier}           & Barrier, native Cholesky \\
\texttt{barrier\_chol\_amd} & \texttt{-chol uni -seconds 28800 -barrier} & Barrier, CHOLMOD + AMD ordering \\
\texttt{primal\_simplex}   & \texttt{-seconds 28800 -primals}           & Primal simplex (default crash) \\
\texttt{dual\_simplex}     & \texttt{-seconds 28800 -duals}             & Dual simplex \\
\bottomrule
\end{tabular}
\end{table}

\noindent\textbf{Important:} In Clp, \texttt{-barrier} and \texttt{-primals}/\texttt{-duals}
are \emph{action} parameters that trigger the solve immediately.  All
configuration parameters (including \texttt{-chol} and \texttt{-seconds})
must therefore precede the action flag.

\subsection{Time Limits}

\begin{itemize}
  \item \textbf{Soft limit:} 28\,800\,s (8\,h), passed to Clp via
        \texttt{-seconds}.
  \item \textbf{Hard limit:} 32\,400\,s (9\,h), enforced externally via
        \texttt{SIGTERM}/\texttt{SIGKILL}.
\end{itemize}

\subsection{Instances}

The eleven instances in Table~\ref{tab:instances} were selected from
MIPLIB~2017 based on structural properties suggested as barrier-friendly
in the \texttt{guess\_option} document: high density, network/flow structure,
high degeneracy, and high equality proportion.  All instances are taken from
\texttt{/home/haroldo/inst/super} in \texttt{.mps.gz} format.

\begin{table}[h]
\centering
\caption{Benchmark instances.}
\label{tab:instances}
\begin{tabular}{@{}lll@{}}
\toprule
\textbf{Instance} & \textbf{Category} & \textbf{Structural rationale} \\
\midrule
\texttt{neos-5075914-elvire}  & Dense / large-scale   & Large, dense constraint matrix \\
\texttt{neos-3083819-nubu}    & Dense / large-scale   & High density, many nz per row \\
\texttt{supportcase18}        & Dense / large-scale   & Large-scale, dense blocks \\
\texttt{graph40-80-1rand}     & Dense / large-scale   & Dense graph-based structure \\
\texttt{mc11}                 & Network / flow        & Multicommodity flow \\
\texttt{momentum1}            & Network / flow        & Large network structure \\
\texttt{neos-4954672-berkel}  & Network / flow        & Flow-like, many equalities \\
\texttt{supportcase6}         & Network / flow        & Large flow/network relaxation \\
\texttt{neos-3627168-kasai}   & Highly degenerate     & Degenerate LP, simplex stalls \\
\texttt{neos-4338804-snowy}   & Highly degenerate     & Degenerate, simplex struggles \\
\texttt{seymour}              & Highly degenerate     & Classic difficult degenerate LP \\
\bottomrule
\end{tabular}
\end{table}

\subsection{Objective Validation}

After each solve that reached optimality, the reported objective value was
compared against a Gurobi LP-relaxation reference
(\texttt{lp\_relaxation\_results\_gurobi.tsv}) using a relative tolerance of
$10^{-4}$.  All optimal solutions in this experiment matched within
$10^{-10}$, confirming correctness across all four configurations.

\subsection{Measurement}

\begin{itemize}
  \item \textbf{Wall-clock time}: measured externally using millisecond
        timestamps around the solver call.
  \item \textbf{Solve time}: Clp-reported time on the
        \texttt{Optimal objective \ldots\ time} output line, excluding
        presolve.
  \item \textbf{Peak memory}: \texttt{Maximum resident set size} from
        \texttt{/usr/bin/time -v} wrapping the Clp process.
\end{itemize}

All 44 jobs (11 instances $\times$ 4 configurations) ran concurrently on a
128-core machine; each Clp process is single-threaded, so at most 44 cores
were active simultaneously.

% ─────────────────────────────────────────────────────────────────────────────
\section{Results}

\subsection{Solve Time}
\label{sec:solvetime}

Table~\ref{tab:solvetime} reports Clp-reported solve time (excluding
presolve).  The fastest configuration for each instance is highlighted in
\colorbox{bestcell}{green}; time-limit hits are in \colorbox{tlcell}{grey}.

\begin{table}[h]
\centering
\caption{Clp-reported solve time (seconds, excluding presolve). TL = time
         limit reached. Best per row highlighted.}
\label{tab:solvetime}
\setlength{\tabcolsep}{5pt}
\begin{tabular}{@{}lrrrrl@{}}
\toprule
\textbf{Instance}
  & \textbf{Barrier native}
  & \textbf{Barrier+AMD}
  & \textbf{Primal}
  & \textbf{Dual}
  & \textbf{Best} \\
\midrule
neos-5075914-elvire  & 0.102 & 0.142 & \cellcolor{bestcell}0.052 & 0.072 & Primal \\
neos-3083819-nubu    & 0.092 & 0.472 & 0.042 & \cellcolor{bestcell}0.032 & Dual \\
supportcase18        & 0.152 & 0.192 & \cellcolor{bestcell}0.042 & 0.152 & Primal \\
graph40-80-1rand     & 25\,452 & \cellcolor{tlcell}TL & \cellcolor{bestcell}14\,148 & \cellcolor{tlcell}TL & Primal \\
mc11                 & \cellcolor{bestcell}0.012 & 0.022 & 0.012 & 0.012 & Barrier native \\
momentum1            & \cellcolor{tlcell}TL & \cellcolor{tlcell}TL & 4.842 & \cellcolor{bestcell}0.602 & Dual \\
neos-4954672-berkel  & 0.022 & 0.182 & \cellcolor{bestcell}0.012 & 0.012 & Primal \\
supportcase6         & 6.742 & 7.582 & \cellcolor{bestcell}1.302 & 3.642 & Primal \\
neos-3627168-kasai   & \cellcolor{bestcell}0.032 & 0.062 & 0.042 & 0.032 & Barrier native \\
neos-4338804-snowy   & 0.612 & 0.432 & \cellcolor{bestcell}0.012 & 0.012 & Primal \\
seymour              & 47.8 & 796 & \cellcolor{bestcell}0.462 & 0.872 & Primal \\
\midrule
\textbf{Wins}        & 2 & 0 & \textbf{7} & 2 & \\
\textbf{Timed out}   & 1 & 2 & 0 & 1 & \\
\textbf{GeoMean ratio} & 4.231 & 11.524 & \cellcolor{bestcell}\textbf{1.270} & 1.388 & \\
\bottomrule
\end{tabular}
\end{table}

\noindent The geometric mean ratio measures each configuration's solve time
relative to the fastest on that instance, computed only over jointly solved
instances.  A value of 1.000 would mean always fastest.

\subsection{Wall-clock Time}
\label{sec:walltime}

Wall-clock times (Table~\ref{tab:walltime}) follow the same pattern but
include startup, I/O, and presolve overhead.  The rankings are essentially
identical to solve time.

\begin{table}[h]
\centering
\caption{Wall-clock time (seconds). TL = time limit reached.}
\label{tab:walltime}
\setlength{\tabcolsep}{5pt}
\begin{tabular}{@{}lrrrrl@{}}
\toprule
\textbf{Instance}
  & \textbf{Barrier native}
  & \textbf{Barrier+AMD}
  & \textbf{Primal}
  & \textbf{Dual}
  & \textbf{Best} \\
\midrule
neos-5075914-elvire  & 0.164 & 0.185 & \cellcolor{bestcell}0.089 & 0.112 & Primal \\
neos-3083819-nubu    & 0.173 & 0.512 & 0.093 & \cellcolor{bestcell}0.076 & Dual \\
supportcase18        & 0.231 & 0.247 & \cellcolor{bestcell}0.099 & 0.178 & Primal \\
graph40-80-1rand     & 25\,529 & \cellcolor{tlcell}TL & \cellcolor{bestcell}14\,219 & \cellcolor{tlcell}TL & Primal \\
mc11                 & 0.047 & 0.051 & \cellcolor{bestcell}0.045 & 0.045 & Primal \\
momentum1            & \cellcolor{tlcell}TL & \cellcolor{tlcell}TL & 5.063 & \cellcolor{bestcell}0.766 & Dual \\
neos-4954672-berkel  & 0.053 & 0.197 & \cellcolor{bestcell}0.036 & 0.053 & Primal \\
supportcase6         & 7.176 & 8.023 & \cellcolor{bestcell}1.743 & 4.110 & Primal \\
neos-3627168-kasai   & 0.081 & 0.099 & 0.116 & \cellcolor{bestcell}0.064 & Dual \\
neos-4338804-snowy   & 0.627 & 0.461 & 0.072 & \cellcolor{bestcell}0.036 & Dual \\
seymour              & 47.842 & 796.738 & \cellcolor{bestcell}0.506 & 0.947 & Primal \\
\bottomrule
\end{tabular}
\end{table}

\subsection{Peak Memory}
\label{sec:memory}

Table~\ref{tab:memory} reports the peak resident set size.  The most
striking differences appear on the three hardest instances
(Table~\ref{tab:memhighlight}).

\begin{table}[h]
\centering
\caption{Peak memory usage (MB). $\uparrow$ highest in row;
         $\downarrow$ lowest in row.}
\label{tab:memory}
\setlength{\tabcolsep}{5pt}
\begin{tabular}{@{}lrrrr@{}}
\toprule
\textbf{Instance}
  & \textbf{Barrier native}
  & \textbf{Barrier+AMD}
  & \textbf{Primal}
  & \textbf{Dual} \\
\midrule
neos-5075914-elvire & 16 & 19$\uparrow$ & 17 & 17 \\
neos-3083819-nubu   & 20 & 19 & 21$\uparrow$ & 21$\uparrow$ \\
supportcase18       & 18 & 23$\uparrow$ & 19 & 18 \\
graph40-80-1rand    & 8\,741 & \cellcolor{badcell}12\,533$\uparrow$ & \textbf{862}$\downarrow$ & 886 \\
mc11                & 17$\uparrow$ & 16 & 17$\uparrow$ & 14$\downarrow$ \\
momentum1           & 786 & \cellcolor{badcell}967$\uparrow$ & \textbf{46}$\downarrow$ & 50 \\
neos-4954672-berkel & 16 & 17$\uparrow$ & 15$\downarrow$ & 17$\uparrow$ \\
supportcase6        & 82 & 81$\downarrow$ & 82 & 84$\uparrow$ \\
neos-3627168-kasai  & 16 & 17$\uparrow$ & 17$\uparrow$ & 14$\downarrow$ \\
neos-4338804-snowy  & 21 & 22$\uparrow$ & 15$\downarrow$ & 16 \\
seymour             & 65 & \cellcolor{badcell}97$\uparrow$ & \textbf{17}$\downarrow$ & 20 \\
\bottomrule
\end{tabular}
\end{table}

\begin{table}[h]
\centering
\caption{Memory overhead of barrier vs.\ primal simplex on hard instances.}
\label{tab:memhighlight}
\begin{tabular}{@{}lrrrr@{}}
\toprule
\textbf{Instance} & \textbf{Barrier native} & \textbf{Barrier+AMD} & \textbf{Primal} & \textbf{AMD/Primal} \\
\midrule
graph40-80-1rand & 8\,741 MB & 12\,533 MB & 862 MB & $\times$14.5 \\
momentum1        & 786 MB    & 967 MB     & 46 MB  & $\times$21.0 \\
seymour          & 65 MB     & 97 MB      & 17 MB  & $\times$5.7  \\
\bottomrule
\end{tabular}
\end{table}

\subsection{Iteration Counts}
\label{sec:iterations}

Table~\ref{tab:iters} reports the number of barrier (IPM) or simplex
iterations.  Barrier's theoretical advantage is its
dimension-independent iteration count.

\begin{table}[h]
\centering
\caption{Number of barrier or simplex iterations. TL = timed out.}
\label{tab:iters}
\setlength{\tabcolsep}{5pt}
\begin{tabular}{@{}lrrrr@{}}
\toprule
\textbf{Instance}
  & \textbf{Barrier native}
  & \textbf{Barrier+AMD}
  & \textbf{Primal}
  & \textbf{Dual} \\
\midrule
neos-5075914-elvire & 20  & 19  & 1\,074  & 1\,376 \\
neos-3083819-nubu   & 446 & 339 & 539     & 528 \\
supportcase18       & 11  & 11  & 663     & 985 \\
graph40-80-1rand    & 11  & TL  & 614\,030 & TL \\
mc11                & 53  & 53  & 662     & 404 \\
momentum1           & TL  & TL  & 7\,293  & 3\,138 \\
neos-4954672-berkel & 60  & 281 & 322     & 409 \\
supportcase6        & 721 & 724 & 9\,787  & 4\,390 \\
neos-3627168-kasai  & 150 & 159 & 1\,294  & 1\,366 \\
neos-4338804-snowy  & 6   & 6   & 474     & 409 \\
seymour             & 49  & 229 & 2\,764  & 3\,710 \\
\bottomrule
\end{tabular}
\end{table}

\noindent The barrier method achieves remarkably few iterations on most
instances (6--53 for the easy ones, 11 for \texttt{graph40-80-1rand}).
However, each iteration involves a Cholesky factorisation that can be orders
of magnitude more expensive per iteration than a simplex pivot, explaining
why low iteration counts do not translate into fast wall-clock times.

% ─────────────────────────────────────────────────────────────────────────────
\section{Analysis}

\subsection{Primal Simplex Dominates}

Primal simplex is the fastest method on 7 of 11 instances and is the only
configuration to solve all 11 within the time limit.  Its geometric mean
ratio to the best (1.27) is far lower than barrier native (4.23) or
barrier+AMD (11.52).  This is consistent with the design of the existing
\texttt{guess()} heuristic, which predominantly recommends primal simplex.

\subsection{Barrier AMD Is Uniformly Worse Than Native Barrier}

On all 9 instances where both barrier variants solved to optimality, native
Cholesky was faster than AMD Cholesky (Table~\ref{tab:native_vs_amd}).
The ratio \texttt{AMD/native} ranges from 1.1$\times$ to 16.7$\times$, with
a median around 1.9$\times$.  Only on \texttt{neos-4338804-snowy} (a trivially
fast instance where both solve in under 1\,s) does AMD show any advantage (0.71$\times$).

\begin{table}[h]
\centering
\caption{Barrier native vs.\ Barrier+AMD: solve time and memory comparison.
         Ratio $>1$ means AMD is slower.}
\label{tab:native_vs_amd}
\setlength{\tabcolsep}{5pt}
\begin{tabular}{@{}lrrrrrr@{}}
\toprule
\textbf{Instance}
  & \textbf{Native (s)}
  & \textbf{AMD (s)}
  & \textbf{Ratio}
  & \textbf{Mem native}
  & \textbf{Mem AMD}
  & \textbf{Mem ratio} \\
\midrule
neos-5075914-elvire  & 0.102 & 0.142 & 1.39$\times$ & 16\,MB & 19\,MB & 1.19$\times$ \\
neos-3083819-nubu    & 0.092 & 0.472 & 5.13$\times$ & 20\,MB & 19\,MB & 0.95$\times$ \\
supportcase18        & 0.152 & 0.192 & 1.26$\times$ & 18\,MB & 23\,MB & 1.28$\times$ \\
graph40-80-1rand     & 25\,452 & TL  & ---          & 8\,741\,MB & 12\,533\,MB & 1.43$\times$ \\
mc11                 & 0.012 & 0.022 & 1.83$\times$ & 17\,MB & 16\,MB & 0.94$\times$ \\
momentum1            & TL    & TL    & ---          & 786\,MB & 967\,MB & 1.23$\times$ \\
neos-4954672-berkel  & 0.022 & 0.182 & 8.27$\times$ & 16\,MB & 17\,MB & 1.06$\times$ \\
supportcase6         & 6.742 & 7.582 & 1.12$\times$ & 82\,MB & 81\,MB & 0.99$\times$ \\
neos-3627168-kasai   & 0.032 & 0.062 & 1.94$\times$ & 16\,MB & 17\,MB & 1.06$\times$ \\
neos-4338804-snowy   & 0.612 & 0.432 & \cellcolor{bestcell}0.71$\times$ & 21\,MB & 22\,MB & 1.05$\times$ \\
seymour              & 47.8  & 796   & \cellcolor{badcell}16.67$\times$ & 65\,MB & 97\,MB & 1.49$\times$ \\
\bottomrule
\end{tabular}
\end{table}

The AMD ordering is designed to minimise fill-in in the Cholesky factor of
the normal-equations matrix $AD^{-2}A^{\top}$.  The results suggest that at
these problem scales the reordering cost exceeds the factorisation savings,
or that the sparsity pattern of the normal equations for these instances is
not amenable to AMD reordering.  The consistently higher memory usage of the
AMD variant (despite its fill-reduction goal) further indicates the
reordering is not producing sparser factors in practice---possibly because
the normal-equations matrix is already dense due to dense columns or
degenerate structure.

\subsection{The \texttt{seymour} Pathology}

\texttt{seymour} is the starkest result.  It is a classic highly-degenerate
covering LP.  Primal simplex solves it in 0.46\,s (2\,764 pivots); native
barrier takes 47.8\,s (49 iterations, 100$\times$ slower); AMD barrier takes
796\,s (229 iterations, 1\,730$\times$ slower).  The AMD variant performs
\emph{four times more barrier iterations} than native and takes 16$\times$
longer per iteration, suggesting the AMD ordering creates a pathologically
ill-conditioned factorisation for this problem's normal-equations matrix.

\subsection{The \texttt{graph40-80-1rand} Case}

This is the instance where barrier comes closest to being competitive.  It
is a dense graph colouring LP with $-512$ as the LP relaxation optimum.
Primal simplex requires 614\,030 pivots and 3.94\,h; native barrier solves
it with only 11 iterations in 7.07\,h.  Barrier's iteration count is
56\,000$\times$ smaller, yet it is 1.8$\times$ slower overall---each
Cholesky factorisation of the dense normal-equations matrix costs enormously.
Native barrier at least terminates; AMD barrier and dual simplex both hit the
8-hour soft limit without converging, and AMD barrier consumes 12.5\,GB of
memory (vs.\ 8.7\,GB native, 862\,MB primal).

This instance illustrates the key trade-off: barrier achieves polynomial
iteration counts on dense problems but each iteration has high
per-iteration cost due to Cholesky, whereas simplex iterates cheaply but may
require an exponential number of pivots in degenerate cases.

\subsection{Network Instances: \texttt{momentum1}}

Both barrier variants fail to converge on \texttt{momentum1} within 8\,hours,
while dual simplex solves it in 0.60\,s.  Network flow problems typically
have a totally unimodular constraint matrix; simplex exploits this structure
via fast pivots with integer intermediate solutions.  Barrier offers no such
structural advantage and the large, sparse network normal-equations matrix
appears difficult to factor efficiently despite its special structure.

\subsection{Memory Overhead}

The memory cost of barrier is significant:
\begin{itemize}
  \item On large instances, barrier uses 10--21$\times$ more memory than
        primal simplex.
  \item AMD Cholesky consistently uses \emph{more} memory than native
        Cholesky, the opposite of its intended purpose, suggesting the
        AMD reordering is not beneficial for these problem structures.
  \item For a MIP solver that calls the LP relaxation at every node,
        barrier's memory overhead would be multiplied across the tree.
\end{itemize}

% ─────────────────────────────────────────────────────────────────────────────
\section{Conclusions and Recommendations}

\subsection{Main Conclusions}

\begin{enumerate}
  \item \textbf{Primal simplex is the best default.}  It is fastest on 7/11
        instances, solves all 11, and has the lowest geometric mean ratio
        (1.27).  The current \texttt{guess()} behaviour of predominantly
        recommending primal simplex is empirically justified.

  \item \textbf{Barrier (native Cholesky) is occasionally competitive.}
        It ties for fastest on \texttt{mc11} and \texttt{neos-3627168-kasai},
        and is the only method to solve \texttt{graph40-80-1rand} within
        the time limit.  However, these gains come with a 10--15$\times$
        memory penalty.

  \item \textbf{AMD Cholesky (\texttt{-chol uni}) provides no benefit over
        native Cholesky} on this benchmark set.  It is slower on 9/9
        comparable pairs and uses more memory in most cases.  It should
        \emph{not} be recommended for these instance types.

  \item \textbf{Barrier is dangerous for degenerate LPs.}  On \texttt{seymour}
        and \texttt{momentum1}, it either diverges slowly or fails to
        converge.  Including barrier in an automated selection strategy
        requires a reliable degeneracy detector.

  \item \textbf{The dead barrier branch in \texttt{guess()} is intentional.}
        Given the evidence, adding a barrier leaf to the decision tree would
        require a strong filter (e.g., dense graph structure detected via the
        OsiFeatures) to avoid catastrophic performance on degenerate or
        network instances.
\end{enumerate}

\subsection{Recommendations for \texttt{guess()} Extension}

If barrier is to be added as a \texttt{guess()} leaf, the following
OsiFeature thresholds are suggested as a filter based on this experiment:
\begin{itemize}
  \item Non-zeros per row $> 50$ (dense matrix);
  \item Fraction of equality rows $> 0.5$;
  \item Fraction of network/flow constraints $< 0.3$ (exclude pure networks);
  \item Degeneracy indicator $< $ threshold (exclude highly degenerate).
\end{itemize}
Even with such a filter, the gain over primal simplex on the one qualifying
instance (\texttt{graph40-80-1rand}) is marginal (barrier is 1.8$\times$
\emph{slower} in wall time despite needing 56\,000$\times$ fewer iterations),
so a barrier leaf may not improve the decision tree's overall performance.

\subsection{AMD Cholesky}

The \texttt{-chol uni} option should be treated as experimental for LP
relaxation solving.  It is not recommended as a default for MIP LP relaxations.
Further investigation on purely continuous (non-MIP) large-scale dense LPs
would be needed to find a regime where it adds value.

% ─────────────────────────────────────────────────────────────────────────────
\section*{Appendix: Raw Data}

\begin{table}[h]
\centering
\caption{Complete results. Status: \texttt{opt} = optimal, \texttt{TL} = soft
         time limit (8\,h), \texttt{HTL} = hard time limit (9\,h).
         Obj.\ match: relative error vs.\ Gurobi reference.}
\label{tab:raw}
\small
\setlength{\tabcolsep}{4pt}
\begin{tabular}{@{}llrrrrrrr@{}}
\toprule
\textbf{Instance} & \textbf{Config} & \textbf{Status} &
\textbf{Wall (s)} & \textbf{Solve (s)} & \textbf{Presolve (s)} &
\textbf{Iters} & \textbf{Mem (MB)} & \textbf{Rel.\ err} \\
\midrule
neos-5075914-elvire & barrier\_default  & opt & 0.164 & 0.102 & 0.01 & 20      & 16  & 1.0e-11 \\
neos-5075914-elvire & barrier\_chol\_amd & opt & 0.185 & 0.142 & 0.01 & 19      & 19  & 1.0e-11 \\
neos-5075914-elvire & primal\_simplex   & opt & 0.089 & 0.052 & 0.01 & 1\,074  & 17  & 1.0e-11 \\
neos-5075914-elvire & dual\_simplex     & opt & 0.112 & 0.072 & 0.01 & 1\,376  & 17  & 1.0e-11 \\
\midrule
neos-3083819-nubu   & barrier\_default  & opt & 0.173 & 0.092 & 0.02 & 446     & 20  & 3.1e-11 \\
neos-3083819-nubu   & barrier\_chol\_amd & opt & 0.512 & 0.472 & 0.02 & 339     & 19  & 3.1e-11 \\
neos-3083819-nubu   & primal\_simplex   & opt & 0.093 & 0.042 & 0.02 & 539     & 21  & 3.1e-11 \\
neos-3083819-nubu   & dual\_simplex     & opt & 0.076 & 0.032 & 0.02 & 528     & 21  & 3.1e-11 \\
\midrule
supportcase18       & barrier\_default  & opt & 0.231 & 0.152 & ---  & 11      & 18  & 7.1e-11 \\
supportcase18       & barrier\_chol\_amd & opt & 0.247 & 0.192 & ---  & 11      & 23  & 7.1e-11 \\
supportcase18       & primal\_simplex   & opt & 0.099 & 0.042 & ---  & 663     & 19  & 7.1e-11 \\
supportcase18       & dual\_simplex     & opt & 0.178 & 0.152 & ---  & 985     & 18  & 7.1e-11 \\
\midrule
graph40-80-1rand    & barrier\_default  & opt & 25\,529 & 25\,452 & 1.97 & 11  & 8\,741 & 0.0e+00 \\
graph40-80-1rand    & barrier\_chol\_amd & TL & 28\,881 & ---   & ---  & ---     & 12\,533 & --- \\
graph40-80-1rand    & primal\_simplex   & opt & 14\,219 & 14\,148 & 2.00 & 614\,030 & 883 & 0.0e+00 \\
graph40-80-1rand    & dual\_simplex     & TL & 28\,863 & ---   & ---  & ---     & 886 & --- \\
\midrule
mc11                & barrier\_default  & opt & 0.047 & 0.012 & 0.00 & 53      & 17  & 3.7e-11 \\
mc11                & barrier\_chol\_amd & opt & 0.051 & 0.022 & 0.00 & 53      & 16  & 3.7e-11 \\
mc11                & primal\_simplex   & opt & 0.045 & 0.012 & 0.00 & 662     & 17  & 3.7e-11 \\
mc11                & dual\_simplex     & opt & 0.045 & 0.012 & 0.00 & 404     & 14  & 3.7e-11 \\
\midrule
momentum1           & barrier\_default  & TL & 29\,185 & ---   & ---  & ---     & 786 & --- \\
momentum1           & barrier\_chol\_amd & TL & 29\,030 & ---   & ---  & ---     & 967 & --- \\
momentum1           & primal\_simplex   & opt & 5.063 & 4.842 & 0.05 & 7\,293  & 46  & 6.2e-11 \\
momentum1           & dual\_simplex     & opt & 0.766 & 0.602 & 0.05 & 3\,138  & 50  & 6.2e-11 \\
\midrule
neos-4954672-berkel & barrier\_default  & opt & 0.053 & 0.022 & 0.00 & 60      & 16  & 2.6e-10 \\
neos-4954672-berkel & barrier\_chol\_amd & opt & 0.197 & 0.182 & 0.01 & 281     & 17  & 2.6e-10 \\
neos-4954672-berkel & primal\_simplex   & opt & 0.036 & 0.012 & 0.00 & 322     & 15  & 2.6e-10 \\
neos-4954672-berkel & dual\_simplex     & opt & 0.053 & 0.012 & 0.00 & 409     & 17  & 2.6e-10 \\
\midrule
supportcase6        & barrier\_default  & opt & 7.176 & 6.742 & ---  & 721     & 82  & 7.7e-11 \\
supportcase6        & barrier\_chol\_amd & opt & 8.023 & 7.582 & ---  & 724     & 81  & 7.7e-11 \\
supportcase6        & primal\_simplex   & opt & 1.743 & 1.302 & ---  & 9\,787  & 82  & 7.7e-11 \\
supportcase6        & dual\_simplex     & opt & 4.110 & 3.642 & 0.28 & 4\,390  & 84  & 7.7e-11 \\
\midrule
neos-3627168-kasai  & barrier\_default  & opt & 0.081 & 0.032 & 0.00 & 150     & 16  & 4.7e-11 \\
neos-3627168-kasai  & barrier\_chol\_amd & opt & 0.099 & 0.062 & 0.00 & 159     & 17  & 4.7e-11 \\
neos-3627168-kasai  & primal\_simplex   & opt & 0.116 & 0.042 & 0.00 & 1\,294  & 17  & 4.7e-11 \\
neos-3627168-kasai  & dual\_simplex     & opt & 0.064 & 0.032 & 0.00 & 1\,366  & 14  & 4.7e-11 \\
\midrule
neos-4338804-snowy  & barrier\_default  & opt & 0.627 & 0.612 & 0.00 & 6       & 21  & 0.0e+00 \\
neos-4338804-snowy  & barrier\_chol\_amd & opt & 0.461 & 0.432 & 0.00 & 6       & 22  & 0.0e+00 \\
neos-4338804-snowy  & primal\_simplex   & opt & 0.072 & 0.012 & 0.00 & 474     & 15  & 0.0e+00 \\
neos-4338804-snowy  & dual\_simplex     & opt & 0.036 & 0.012 & 0.00 & 409     & 16  & 0.0e+00 \\
\midrule
seymour             & barrier\_default  & opt & 47.842 & 47.772 & 0.01 & 49    & 65  & 6.2e-11 \\
seymour             & barrier\_chol\_amd & opt & 796.738 & 796.382 & 0.02 & 229 & 97  & 6.2e-11 \\
seymour             & primal\_simplex   & opt & 0.506 & 0.462 & 0.01 & 2\,764  & 17  & 6.2e-11 \\
seymour             & dual\_simplex     & opt & 0.947 & 0.872 & 0.01 & 3\,710  & 20  & 6.2e-11 \\
\bottomrule
\end{tabular}
\end{table}

\end{document}
